\section{INTRODUCTION}
\label{sec:intro}

Sparse inputs provide limited information about visible regions, causing 3D reconstruction methods to overfit local observations and struggle to maintain global geometric consistency~\cite{Yang23_FreeNeRF}. Sparse observations also introduce greater geometric ambiguity, making occluded regions, textureless areas, and view boundaries prone to missing geometry, structural distortions, and degraded appearance~\cite{Niemeyer22_RegNeRF}. \rev{Kerbl et al.}~\cite{Kerbl23_3DGS} \rev{proposed 3D Gaussian Splatting (3DGS), which achieves fast training and real-time rendering through a tile-based rasterization mechanism. Subsequent developments have further enhanced the representation through structured Gaussian organization}~\cite{Lu24_ScaffoldGS}. Under sparse views, Structure-from-Motion (SfM)~\cite{Schonberger16_COLMAP} often recovers incomplete point clouds. Since 3DGS relies on SfM cameras and sparse points for initialization, weak coverage harms geometric convergence~\cite{Chung24_DepthRegGS}. Limited cross-view constraints further cause floaters and pseudo-geometry~\cite{Xu25_DropoutGS}, while sparse supervision~\cite{Park25_DropGaussian} and visibility imbalance aggravate overfitting~\cite{Somraj23_ViPNeRF}.

Existing sparse-view methods have explored depth regularization~\cite{Li24_DNGaussian}, Gaussian densification~\cite{Zhu24_FSGS}, co-regularization~\cite{Zhang24_CoRGS}, and Gaussian dropout~\cite{Park25_DropGaussian} to improve reconstruction. Nevertheless, incomplete initial geometry and overfitting during optimization remain challenges when only a few observations are available. \rev{To mitigate these challenges, we develop SynGS within the sparse-view 3DGS setting, where geometry-guided Gaussian initialization is coupled with regularization during training.} By strengthening initial geometric support and mitigating overfitting during optimization, SynGS improves reconstruction robustness under severely limited view coverage.
The main contributions are:
\begin{itemize}
\item \textbf{Visual Hull Initialization.} We incorporate a visual hull derived from multi-view silhouettes to provide an explicit geometric constraint for Gaussian initialization. This constraint complements incomplete SfM observations and helps establish a more reliable initial spatial distribution of Gaussian primitives.
\item \textbf{VGGT-Guided Densification.} We develop a pre-optimization densification strategy that back-projects cross-view-consistent depth estimates from VGGT to supplement the sparse initial point set. By introducing additional geometric samples before training, this strategy improves geometric coverage in under-observed regions.
\item \textbf{Dynamic Visibility Regularization.} We incorporate Dynamic Visibility Regularization during training to stochastically mask Gaussian primitives, redistribute optimization gradients, and mitigate overfitting under limited-view supervision. \rev{Together with the geometry-enhanced initialization, this regularization contributes to more robust reconstruction when only a limited number of input views are available.}
\end{itemize}
